\ifx\danslelivre\undefined
\documentclass[10pt]{../fichiers-configuration-latex/smfart}
\usepackage{ifthen}
\newboolean{english}
\newboolean{francais}
\setboolean{english}{false}
\setboolean{francais}{true}
\usepackage{palatino,euler}
\input{../fichiers-configuration-latex/paquets}
\selectlanguage{francais}
\input{../fichiers-configuration-latex/lemme-proposition-theoreme}
\input{../fichiers-configuration-latex/numerotation}
\input{../fichiers-configuration-latex/macros}
\input{../fichiers-configuration-latex/ucs_manquants}
\input{../.cv}


\externaldocument{uniformisation-faible}
\externaldocument{introduction}

\begin{document}

\begin{center}
IV. Le théorème de Cohen-Gabber\\
Fabrice Orgogozo
\end{center}

\version
\tableofcontents

\else
\setboolean{english}{false}
\setboolean{francais}{true}
\selectlanguage{francais}
\chapter{Le théorème de Cohen-Gabber}
\addtocontents{toc}{par Fabrice Orgogozo}
\begin{center}
Fabrice Orgogozo
\end{center}
\fi

\section{$p$-bases et différentielles (rappels)}

\subsection{Définition et caractérisation différentielle}

\subsubsection{}
Pour la commodité du lecteur, et pour fixer les
notations, nous rappelons ici quelques résultats bien
connus dont nous ferons usage ci-après.
Nous conseillons au lecteur de ne s'y reporter qu'en
cas de besoin.

\begin{définition2}
Soient $k$ un corps de caractéristique $p>0$,
$K$ une extension de $k$, et $(b_i)_{i∈I}$ une famille d'éléments de $K$.
On dit que les $(b_i)$ constituent une
\emph{$p$-base} de $K$ sur $k$ (resp. sont \emph{$p$-libres} sur $k$) 
si les monômes $\prod_i b_i^{n(i)}$ ($0≤n(i)<p$,
$(n(i))_{i∈I}$ de support fini) forment une base
du $k(K^p)$-espace vectoriel $K$ (resp. sont linéairement indépendants
sur $k(K^p)$).
\end{définition2}

Si $k=\FF_p$, on parle alors de $p$-base \emph{absolue},
ou de $p$-base s'il n'y a pas d'ambiguïté.
Enfin, on appelle parfois \emph{$p$-monôme} un produit comme
ci-dessus. Un lien entre cette notion et la structure des anneaux
locaux complets ressort du théorème suivant.

\begin{theoreme2}[\ac{v}{3}{3}{, th. 1 b)}]
Soient $A$ un anneau local séparé complet de caractéristique $p>0$
et $(β_i)_{i∈I}$ une famille d'éléments de $A$ dont les classes modulo
l'idéal maximal $𝔪 _A$ forment une $p$-base du corps résiduel $A ∕ 𝔪_A$.  Il existe
alors un unique corps de représentants de $A$ contenant les éléments $β_i$.
\end{theoreme2}

\begin{remarque2}
On peut étendre de façon évidente la notion de $p$-base
au cas d'un anneau quelconque de caractéristique $p>0$,
cf. \ega{$0_{\textsc{iv}}$}{21.1-4}. Nous n'en aurons pas besoin.
\end{remarque2}

\subsubsection{}On vérifie immédiatement que les $(b_i)_{i∈I}$
forment une $p$-base de $K$ sur $k$ si et seulement si, pour tout $i∈I$,
l'élément $b_i$ n'appartient pas au sous-corps $k(K^p,(b_j)_{j≠i})$
de $K$. 
(Voir p. ex. \ega{$0_{\textsc{iv}}$}{21.4.3}.)

\subsubsection{}Pour toute extension de corps $K/k$, nous noterons $d_{K/k}$ la différentielle $K→ Ω¹_{K/k}$.

\begin{proposition2}
\label{4-caractérisation différentielle}
Soient $k$ un corps de caractéristique $p>0$, et $K$ une extension de $k$.
Une famille $(b_i)_{i∈I}$ d'éléments de $K$
est une $p$-base de $K$ sur $k$ si et seulement si les différentielles
$d_{K/k}(b_i)$ forment une base du $K$-espace vectoriel $Ω¹_{K/k}$.
\end{proposition2}

\begin{démo}
Soit $B=(b_i)_{i∈I}$ une $p$-base de $K$ sur $k$. Tout morphisme
(ensembliste) $Δ:B→K$ s'étend de manière unique en une $k$-dérivation
$D$ de $K$ : il suffit de poser $D(b₁^{n₁}\cdots b_r^{n_r})= ∑_i n_i
b₁^{n₁}\cdots b_i^{n_i-1}\cdots b_r^{n_r}$ et de 
l'étendre par $k(K^p)$-linéarité. Cela est équivalent au fait que les $d_{K/k}(b_i)$
forment une base de $Ω¹_{K/k}$. (Le fait que $K$ soit un corps
n'est pas utilisé dans cette implication.) Réciproquement, si les $d_{K/k}(b_i)$ forment une base, on 
observe que les $p$-monômes sont $k(K^p)$-linéairement indépendants : 
dans le cas contraire on aurait, pour un indice $i$ convenable,
$b_i∈K^p((b_j)_{j≠i})$, ce qui se traduirait par une relation
linéaire entre les différentielles. Soit $B'$ une $p$-base de $K$ sur $k$ contenant les
$b_i$ (\emph{loc. cit.}, 21.4.2) ; d'après l'implication précédente, on a nécessairement
$B'=(b_i)_{i∈I}$.
\end{démo}

\begin{corollaire2}
\label{4-dérivée nulle}
Soient $k$ un corps de caractéristique $p>0$, $K$ une extension de $k$.
Un élément $x$ de $K$ appartient à $k(K^p)$ si et seulement si
$d_{K/k}(x)=0$.
\end{corollaire2}

\subsubsection{}Rappelons que le \emph{$p$-rang} d'un corps est le cardinal
d'une $p$-base absolue (bien défini en vertu de ce qui précède).
On vérifie immédiatement que ce \emph{cardinal} (fini ou non) est invariant par
extension \emph{finie} de corps.

\subsection{Stabilisation}

\begin{lemme2}[cf. p. ex. \ega{$0_{\textsc{iv}}$}{21.8.1}]
Soient $K$ un corps, $k$ un sous-corps, $(k_α)_{α∈I}$ une famille
de sous-corps de $K$ telle que $⋂k_α=k$ et
\emph{filtrante décroissante}, 
c'est-à-dire telle que pour toute paire d'indices $α,β$, il existe un indice
$γ$ tel que $k_γ ⊂ k_α ∩ k_β$.
Soient $V$ un $K$-espace
vectoriel,  et $(v_i)$ ($1≤i≤n$) une famille finie de vecteurs de $V$.
Si la famille $(v_i)$ est libre sur $k$, il existe un indice $γ$ telle qu'elle
soit aussi libre sur $k_γ$.
\end{lemme2}

\begin{lemme2}
Soient $K$ un corps de caractéristique $p>0$, $k$ un sous-corps et 
$(K_α)_{α∈I}$ une famille filtrante décroissante de sous-corps contenants
$k$.
Les conditions suivantes sont équivalentes :
\begin{enumerate}
\item $⋂_α K_α(K^p)=k(K^p)$ ; 
\item pour tout ensemble fini $\{b₁,\dots,b_n\}⊂K$, $p$-libre sur $k$, il existe
un indice $α$ tel qu'il soit $p$-libre sur $K_α$ ;  
\item il existe une $p$-base de $K$ sur $k$ telle que tout sous-ensemble fini 
soit $p$-libre sur un $k_α$ pour $α$ convenable ; 
\item le morphisme canonique $Ω¹_{K/k}→ \lim_α   Ω¹_{K/K_α}$ est \emph{injectif}.
\end{enumerate}
\end{lemme2}

\begin{démo}
(i)⇒(ii) est une conséquence immédiate du lemme précédent. (ii)⇒(iii) est
trivial (toute $p$-base convient). (iii)⇒(iv) trivial (utiliser
\ref{4-caractérisation différentielle}). Vérifions (iv)⇒(i). Soit $x\notin
k(K^p)$. D'après \ref{4-dérivée nulle}, $d_{K/k}(x)≠0$ de sorte qu'il existe
$α$ tel que $d_{K/k_α}(x)$ soit également non nul. D'après \emph{loc. cit.},
cela entraîne que $x\notin k_α(K^p)$.
\end{démo}

On en déduit le lemme suivant, qui est un cas particulier
de \egazero{iv}{21.8.5}.

\begin{lemme2}
Soient $K$ un corps de caractéristique $p$, $k$ un sous-corps
 et $(K_α)_{α∈I}$ une famille filtrante décroissante de
sous-corps de $K$ contenant $k$ telle que $⋂_α K_α(K^p)=k(K^p)$. Pour toute extension finie
$L$ de $K$, on a également \mbox{$⋂_α K_α(L^p)=k(L^p)$}.
\end{lemme2}

\begin{démo}
On se ramène immédiatement au cas où $L/K$ est monogène.
Si $L/K$ est (algébrique) séparable, la conclusion résulte immédiatement
de l'existence des isomorphismes canoniques $Ω¹_{L/k}\iso Ω¹_{K/k}⊗_K L$,
$Ω¹_{L/K_α}\iso Ω¹_{K/K_α}⊗_K L$ et du critère (iv) ci-dessus.
Dans le cas contraire, $L=K(a)$, où $b=a^p∈K-K^p$.
On distingue naturellement deux cas.
Premier cas : $d_{K/k}(b)=0$, c'est-à-dire $b ∈ k(K^p)$.
Il en résulte que pour toute sous-extension $M$ de $L ∕ k$, on a l'égalité
$M(L^p)=M(K^p,b)=M(K^p)$. Ainsi,
\[⋂_α K_α(L^p)=⋂_α K_α(K^p)=k(K^p)=k(L^p),\]
et la conclusion résulte du critère (i) ci-dessus.
Second cas : $d_{K/k}(b)≠0$. On peut alors compléter $\{b\}$ en une $p$-base
de $K$ sur $k$, que l'on note $(b,(b_i)_{i∈I})$. La famille
$(a,(b_i)_{i∈I})$ est alors une $p$-base de $L$ sur $k$
et on vérifie immédiatement le critère (iii) ci-dessus :
si $(b,b₁,\dots,b_n)$ est $p$-libre sur $k_α$, il en est de même
de $(a,b₁,\dots,b_n)$.
\end{démo}

\begin{proposition2}[\cite{KKK@Matsumura}, §30, lemme 6]\label{4-lemme Nagata}
Soient $K$ un corps de caractéristique $p>0$ et
$(K_α)$ une famille filtrante décroissante de sous-corps
\emph{cofinis} — c'est-à-dire tels que les degrés $[K:K_α]$
soient finis — telle que $⋂_α K_α=K^p$.
Alors, pour toute extension finie $L/K$, il existe un
indice $β$ tel que pour tout sous-corps cofini $K′⊂K_{β}$ on ait :
\[
\rang_L Ω¹_{L/K′}=\rang_{K} Ω¹_{K/K′}.
\]

\end{proposition2}

\begin{démo}
On souhaite se ramener au cas où $L$ est monogène sur $K$.
Pour cela considérons une sous-$K$-extension $M$ de $L$ et posons $M_α=K_α(M^p)⊂M$.
Ce sont des sous-corps cofinis de $M$ et, pour $α,β$ et $γ$ comme
dans l'énoncé, on a $M_γ⊂M_α∩M_β$. En vertu du lemme
précédent, appliqué dans le cas particulier où $k=K^p$,
on a l'égalité $⋂_α M_α=M^p$. On se ramène au cas où $L/M$ est monogène
(par récurrence) en remarquant également que les extensions $M_α/K_α$ sont
\emph{finies}, et que pour tout sous-corps $K′$ de $K$ on a
$Ω¹_{M/K′}=Ω¹_{M/K′(M^p)}$.

Supposons dorénavant l'extension $L ∕ K$ monogène.
Si elle est (algébrique) séparable, le théorème est trivial :
on a $Ω¹_{L/K′} ⭉  Ω¹_{K/K′}⊗_K L$ pour tout $K′⊆ K$.
Sinon, $L=K(a)$, où $a^p=b∈K-K^p$,
et, pour chaque $K′ ⊆ K$, $Ω¹_{L ∕ K ′}$ est naturellement isomorphe à
\[
\big(Ω¹_{K ∕ K′}/Kd_{K ∕ K′}(b)\big)\! ⊗_K \! L   ⊕ L d_{L ∕ K′}(a).
\]
De plus, $d_{K ∕ K^p}(b)$ et $d_{L ∕ L^p}(a)$ sont non nuls
car $b$ (resp. $a$) n'appartient pas à $K^p$ (resp. $L^p$).
Puisque $K^p=⋂K_α$ (resp. $L^p=⋂ K_α(L^p)$) , il existe
un $β$ tel que $d_{K ∕ K_β}(b) ≠ 0$ (resp. $d_{L ∕ K_β}(a)≠
0$). Il résulte de l'isomorphisme ci-dessus que pour chaque $K′ ⊆ K_β$, on a
l'égalité $\rang_L Ω¹_{L ∕ K ′}=\rang_K Ω¹_{K ∕ K ′}$. CQFD.
\end{démo}

\subsubsection{}\label{4-dérivations continues}
Rappelons enfin que si $A$ et $B$ sont deux anneaux linéairement topologisés
(\egazero{i}{7.1.1}), et $A→B$ un morphisme \emph{continu}, 
le $B$-module $Ω¹_{B/A}$ est un $B$-module
\emph{topologique}, la topologie étant déduite 
de celle de $B⊗_A B$ par restriction et passage au quotient 
(\egazero{iv}{20.4.3}). Le $B$-module sous-jacent ne dépend pas des topologies de $A$ et $B$.
On note $\chap{Ω}¹_{B/A}$ son \emph{séparé complété} ; 
il est isomorphe à une limite de $Ω¹$ de morphismes
entre anneaux topologiques 
discrets (\emph{loc. cit.}, 20.7.4).
Pour tout $B$-module topologique $L$, le morphisme canonique
induit par la dérivation universelle est un \emph{isomorphisme} :
$$
\mathrm{Hom.cont}_{B}(Ω¹_{B/A},L)\iso \mathrm{D\acute{e}r.cont}_{A}(B,L).
$$
Si $B$ et $L$ sont séparés et complets, le terme de gauche s'identifie
canoniquement à $\mathrm{Hom.cont}_{B}(\chap{Ω}¹_{B/A},L)$.
Comme on le constate dans le cas particulier très simple où $A$ est un corps
et $B$ un anneau de séries formelles, le $B$-module $\chap{Ω}¹_{B/A}$ a des propriétés de finitude
bien plus remarquables que $Ω¹_{B/A}$ (\emph{loc. cit.}, exemple 20.7.6 et prop.
20.7.5).

\section{Les théorèmes de Cohen-Gabber en caractéristique $>0$}

\subsection{Le théorème de Cohen-Gabber non équivariant en caractéristique $>0$}
Le but de ce paragraphe est de démontrer la variante suivante du théorème de structure
des anneaux locaux nœthériens complets %(\ac{ix}{3}{3}{.th. 2 b)}
\egazero{iv}{19.8.8} (ii), dû à Irving S. Cohen.

\begin{theoreme2}[théorème de Cohen-Gabber ;  \cite{Conference-Illusie@Gabber}, lemme 8.1]
\label{4-Cohen-Gabber égale}
Soit $A$ un anneau local complet nœthérien \emph{réduit}, d'égale
caractéristique $p>0$, équidimensionnel
de dimension $d$ et de corps résiduel $k$.
Il existe un sous-anneau $A₀$ de $A$, isomorphe à $k[[t₁,\dots,t_d]]$,
tel que $A$ soit fini sur $A₀$, sans torsion et génériquement étale.
De plus, le morphisme $A₀→A$ induit un isomorphisme sur les corps résiduels.
\end{theoreme2}

\begin{remarques2}Ce résultat apparaît explicitement comme hypothèse,
pour $A$ intègre, dans \ega{$0_{\textsc{iv}}$}{21.9.5}.
L'expression « génériquement étale » signifie ici qu'il existe un
ouvert dense de $\Spec(A₀)$ au-dessus duquel le morphisme
$\Spec(A)→\Spec(A₀)$ est étale.
\end{remarques2}

\subsubsection{}La démonstration du théorème, qui est une adaptation au cas non irréductible
de \cite{Conference-Illusie@Gabber}, occupe le reste de cette section.
Nous supposerons par la suite $d>0$, sans quoi l'énoncé est évident.
Dans les alinéas \ref{4-début Cohen-Gabber} à \ref{4-fin Cohen-Gabber},
nous allons montrer qu'il existe un corps de représentants $κ$ de $A$ tel
que le $A$-module des formes différentielles
complété $\chap{Ω}¹_{A/κ}$ soit de rang générique
égal à $d$ sur chaque composante irréductible.
En (\ref{4-fin démo}) nous verrons comment en déduire
rapidement le théorème.

\subsubsection{}\label{4-début Cohen-Gabber}Soit $(b_i)_{i∈E}$ une $p$-base de $k=A/\MM_A$.
Choisissons des relèvements arbitraires $β_i$
des $b_i$ dans $A$. Rappelons qu'il existe un unique corps de
représentants $κ\subset A$ contenant les $β_i$ et se surjectant sur $k$
(cf. \ac{ix}{2}{2}{, th. 1 a)}). Changer de corps de
représentants revient donc à changer les $β_i$.
Fixons également un système de paramètres $τ₁,\dots,τ_d$
de $A$ ;  nous ne le changerons qu'à la fin de la
démonstration (\ref{4-fin démo}).

\subsubsection{}Pour toute partie \emph{finie} $e\subset E$, posons $κ_e:=κ^p(β_i,\,i\notin e)\subset κ$.
Les trois propriétés suivantes sont évidentes :
$$\text{pour toute partie finie $e⊂E$}, [κ:κ_e]<+∞,$$
$$\text{pour toutes parties finies $e,e'⊂E$}, κ_{e\cup
e'} ⊂ κ_{e}\cap κ_{e'},$$
$$⋂_{e\subset E} κ_e = κ^p.$$

\subsubsection{}Soient $\Spec(\sur{A})$ une composante
irréductible de $\Spec(A)$, munie de la structure
réduite, et $\sur{τ₁},\dots,\sur{τ_d}$ les images des $τ_i$
dans $\sur{A}$ par la surjection canonique $A↠\sur{A}$.
Considérons le diagramme d'anneaux :
\[
\xymatrix{
\kappa_e[[\sur{\tau_1}^p,\dots,\sur{\tau_d}^p]] \ar[r] \ar[d] &
\kappa[[\sur{\tau_1},\dots ,\sur{\tau_d}]]  \ar[r]
\ar[d] & \sur{A} \ar[d] \\
L_{\kappa,e} \ar[r] & L_\kappa \ar[r] & L 
}
\]
où les flèches horizontales sont les homomorphismes
canoniques, et les flèches verticales les inclusions
dans les corps de fractions respectifs. Les flèches
horizontales sont injectives et correspondent à
des morphismes finis. Pour la seconde, cela résulte
du fait que le module $\sur{A}$ est \emph{quasi-fini}
(\ega{$0_{\textrm{I}}$}{7.4.1}) sur $κ[[\sur{τ₁},\dots
,\sur{τ_d}]]$ donc de type fini car l'idéal $(\sur{τ₁},\dots
,\sur{τ_d})\sur{A}$ est un idéal de définition
(\ega{$0_{\textrm{I}}$}{7.4.4}). Enfin, 
les $\sur{τ_i}$ sont analytiquement indépendants sur
$κ$ : le sous-anneau $κ[[\sur{τ₁},\dots
,\sur{τ_d}]]$ de $\sur{A}$ est bien un anneau de séries
formelles (\ega{$0_{\textrm{IV}}$}{16.3.10}).

On a observé ci-dessus que la famille des $κ_e⊂κ$, $e⊂E$,
satisfait aux hypothèses de la proposition \ref{4-lemme
Nagata}. On vérifie
immédiatement qu'il en est de même de la famille
des sous-corps $L_{κ,e}$ de $L_κ$ ;  on a donc l'égalité
\begin{equation}\label{4-equation1}
\rang_L Ω¹_{L/L_{κ,e}}=\rang_{L_κ} Ω¹_{L_κ/L_{κ,e}},
\end{equation}
dès que l'ensemble fini $e$ est suffisamment grand.

Posons $R_κ=\kappa[[\sur{τ_1},\dots
,\sur{τ_d}]]$ et
$R_{κ,e}=\kappa_e[[\sur{\tau_1}^p,\dots,\sur{\tau_d}^p]]$.
Le terme de gauche de (\ref{4-equation1})
est le rang générique du $\sur{A}$-module
$Ω¹_{\sur{A}/R_{κ,e}}$, c'est-à-dire le rang de son tensorisé
avec $L$.
Remarquons que d'après \ega{$0_{\textrm{IV}}$}{21.9.4}, 
$Ω¹_{\sur{A}/R_{κ,e}}$ s'identifie au module $\sur{A}$-module
$\chap{Ω}¹_{\sur{A}/κ_e}$ de formes différentielles complété.
Le terme de droite est quant à lui le 
rang du $R_{κ,e}$-module libre $Ω¹_{R_{κ,e}/R_κ}$.
Ce dernier est égal à $d+\rang_κ Ω¹_{κ/κ_e}=d+|e|$
(où $|-|$ désigne le cardinal d'un ensemble), de sorte que 
la formule \ref{4-equation1} se réécrit :
\begin{equation}\label{4-equation2}
\rang_{\sur{A}} \chap{Ω}¹_{\sur{A}/κ_e}=d+|e|.
\end{equation}


\subsubsection{}La proposition suivante va nous permettre de modifier
le corps des représentants de façon à pouvoir supposer
$e$ vide (de façon équivalente : $κ_e=κ$).

\begin{proposition2}\label{4-changement relèvements}
Il existe une partie finie $e$ de $E$ et des éléments
$β'_i$, pour $i∈e$, relevant les $b_i$ tels que,
pour chaque composante irréductible intègre $\Spec(\sur{A})$ de
$\Spec(A)$,
les conditions suivantes soient vérifiées :
\begin{enumerate}
\item $\rang_{\sur{A}} \chap{Ω}¹_{\sur{A}/κ_e}=d+|e|$,
\item les images des $dβ'_i$ dans
$\chap{Ω}¹_{\sur{A}/κ_e}⊗_{\sur{A}}
L$, où $L=\Frac(\sur{A})$, sont $L$-linéairement indépendantes.
\end{enumerate}
\end{proposition2}

L'égalité \ref{4-equation1} (et donc \ref{4-equation2}) étant valable, pour chaque
composante irréductible, dès que $e$ est suffisamment 
grand, on peut choisir un tel ensemble qui convient
pour chacune d'entre elles.
La propriété (i) en découle.

Pour démontrer la propriété (ii), nous utiliserons le lemme élémentaire suivant.

\begin{lemme2}
Soient $\sur{A}$ et $L$ comme ci-dessus. 
Pour tout idéal non nul $I$ de $\sur{A}$, l'ensemble
des $df\otimes_{\sur{A}} L$, pour $f∈I$, est une famille 
génératrice du $L$-espace vectoriel
$\chap{Ω}¹_{\sur{A}/κ_e}\otimes_{\sur{A}} L$.
\end{lemme2}

\begin{proof}
Soient $f₀∈I$ non nul, et $ω₀=df₀$. Pour tout 
$b∈\sur{A}$, $d(bf₀)=bω₀+f₀db$.
La famille des $d(bf₀)⊗1$ 
contient $ω₀⊗1$ ;  d'après la formule précédente,
le $L$-espace vectoriel qu'elle engendre
contient donc les $db⊗1$ pour chaque $b∈\sur{A}$.
\end{proof}

Soit $\{𝔭₁,\dots,𝔭_c\}$ l'ensemble des idéaux premiers
minimaux de $A$. 
Pour chaque $j∈\{1,\cdots,c\}$, posons $A_j=A/𝔭_j$ et
$X_j=\Spec(A_j)$
la composante irréductible intègre de $X=\Spec(A)$ correspondante.
Notons pour tous $i∈e$ et $j∈\{1,\dots,c\}$,
$β_{i,j}$ l'image dans $A_j$ de $β_i∈A$.
(Rappelons que les $β_i$ font partie d'une $p$-base de $κ\subset A$.)
Nous allons démontrer par récurrence sur $j$ ($0≤j≤c$) qu'il existe des éléments 
$\{m_{i,j}\}$ dans $\MM_A$, pour $i∈e$, tels que
les images des éléments $β_i+m_{i,j}$ dans chacun des
anneaux $A_1,\dots,A_j$ aient des différentielles
linéairement indépendantes dans chacun des
espaces vectoriels
$Ω¹_{A₁/R_{κ,e}}\otimes_{A₁} \Frac{\,A₁}$,
\dots, $Ω¹_{A_j/R_{κ,e}}\otimes_{A_j} \Frac{\,A_j}$. 
Pour $j=0$, cette condition est vide.
Supposons l'assertion démontrée pour un $j≤c-1$ et
montrons la pour $j+1$.
Quitte à remplacer $β_i$ par $β_i+m_{i,j}$, on peut 
supposer que $m_{i,j}=0$ pour tout $i∈e$.
L'anneau $A$ étant réduit, les $𝔭_α$ forment une
décomposition primaire \emph{réduite}
de $(0)$, de
sorte que l'idéal $\got{q}_j:=𝔭_1\cap \cdots \cap 𝔭_j$
($=\mathrm{Ker}(A→A_1×\cdots×A_j)$) 
n'est pas contenu dans $𝔭_{j+1}$.
Si $j>0$, notons $I_{j+1}$ son image dans
$\sur{A}=A_{j+1}(=A/𝔭_{j+1})$ ;  c'est un idéal non nul. Si $j=0$,
on considère $\MM_{\sur{A}}$. 
D'après (i), $\rang_{\sur{A}} \chap{Ω}¹_{\sur{A}/κ_e}=d+|e|≥|e|$ ; 
d'autre part, la famille $d(I_{j+1})$ est génératrice
dans $\chap{Ω}¹_{\sur{A}/κ_e}\otimes_{\sur{A}} L$
(où $L=\Frac{\,\sur{A}}$). 

\begin{lemme2}
Soient $V$ un espace vectoriel de dimension au moins
$n$, $b₁,\dots,b_n$ des vecteurs de $V$ et $W$ une
famille génératrice. Il existe une famille
$w₁,\dots,w_n$ d'éléments de $W ∪ \{0\}$ tels que 
les $b_i+w_i$ soient linéairement indépendants.
\end{lemme2}

\begin{démo}Par récurrence immédiate sur $n$.\end{démo}

Il existe donc des éléments $m'_{i,j+1}∈I_{j+1}$, $i∈e$,
tels que les différentielles des éléments
\mbox{$d((β_i \mod 𝔭_{j+1})+m'_{i,j+1})$}, $i∈e$, soient linéairement indépendantes
dans $\chap{Ω}¹_{\sur{A}/κ_e}\otimes_{\sur{A}} L$.

Relevons les $m'_{i,j+1}$ en des éléments $m_{i,j+1}$
de $\got{q}_j$ si $j>0$, ou de $\MM_A$ si $j=0$. Par construction,
ils satisfont la propriété escomptée au cran $j+1$.


\subsubsection{}\label{4-fin Cohen-Gabber}
Considérons le sous-corps $κ':=κ^p(β_i, i∉e;β'_i, i∈e)=κ_e(β'_i, i∈e)\subset A$, où les $β'_i$ ($i∈e$) sont
comme en \ref{4-changement relèvements}. Il s'envoie isomorphiquement sur
$k=A/\MM_A$ par réduction : son image contient $k^p$ et
les images des $β_i$ ($i∉e$), $β'_i$ ($i∈e$), qui constituent une $p$-base
de $k$.
Des égalités \ref{4-equation2} et de la propriété (ii) de 
\ref{4-changement relèvements},  on tire : 
$$\rang_{\sur{A}} \chap{Ω}¹_{\sur{A}/κ'}=d,$$ pour toute composante irréductible
intègre
$\Spec(\sur{A})$ de $X$. 
Par la suite, nous noterons encore $κ$ ce nouveau
corps de représentants.

\subsubsection{}\label{4-fin démo}Le $A$-module $\chap{Ω}¹_{A/κ}$ étant de rang générique $d$ sur chaque
composante irréductible, on montre en procédant comme précédemment, 
qu'il existe des éléments $f_1,\dots,f_d$ de $A$ tels que les $d(f_i \ \mathrm{mod}
\, 𝔭_α)\otimes_{A_j} \Frac{\,A_j}$ forment une base de 
$\chap{Ω}¹_{A_j/κ} \otimes_{A_j} \Frac{\,A_j}$ pour chaque composante
irréductible $\Spec(A_j)$ de $X$. Quitte à les
multiplier individuellement par une puissance $p$-ième
d'un élément appartenant à $\MM_A-⋃_j 𝔭_j$, on peut les supposer dans $\MM_A$.
Rappelons que l'on a choisi un système de paramètres
$τ_1,\dots,τ_d$ dans $A$, de sorte que le morphisme
$\Spec(A)→\Spec(k[[τ_1,\dots,τ_d]])$ soit fini.

Posons, pour $i∈\{1,\dots,d\}$,
$$
t_i:=τ_i^p(1+f_i).
$$
Soient $A₀$ le sous-anneau $κ[[t_1,\dots,t_d]]$ de
$A$, $X₀=\Spec(A₀)$. Le morphisme $X→X₀$ est fini : cela résulte du 
fait que les éléments $1+f_i$ sont des unités de $A$. Vérifions 
qu'il est génériquement étale.
L'anneau $A$ étant nœthérien complet, le $A$-module de type
fini $Ω¹_{A/A₀}$ est également complet et coïncide donc
avec le module des formes différentielles complété $\chap{Ω}¹_{A/A₀}$.
Les anneaux $A₀$ et $A$ étant métrisables, et tout sous-$A$-module
de $\chap{Ω}¹_{A/κ}$ étant fermé, la suite $$
\chap{Ω}¹_{A₀/κ}\chap{⊗}_{A₀} A →  \chap{Ω}¹_{A/κ} →
\chap{Ω}¹_{A/A₀}=Ω¹_{A/A₀}→0
$$
est \emph{exacte} (\ega{$0_{\textrm{IV}}$}{20.7.17}).
Il résulte de l'hypothèse sur les éléments $f_i$ et de la formule
$$
d(t_i)=τ_i^p df_i
$$
qu'au-dessus de chaque point maximal de $X=\Spec(A)$, la première flèche
est surjective. On en déduit que le $A$-module $Ω¹_{A/A₀}$ est génériquement nul,
CQFD.

\subsection{Le théorème de Cohen-Gabber équivariant en caractéristique $>0$}

\subsubsection{}Nous allons démontrer ici une généralisation
du théorème \ref{4-Cohen-Gabber égale} dans le cas d'un anneau non nécessairement
équidimensionnel, muni d'une action d'un groupe fini.

\begin{theoreme2}\label{4-Cohen-Gabber équivariant}
Soient $A$ un anneau local nœthérien complet réduit, de dimension $d$, de corps résiduel $κ$ et 
$G$ un groupe fini agissant sur $A$ avec $|G|$ inversible dans $κ$. Alors, il
existe un morphisme fini génériquement étale, $G$-équivariant, $κ[[t₁,\dots,t_d]]→ A$, où $κ→A$ relève l'identité
de $κ$ et $G$ agit trivialement sur les $t_i$.
\end{theoreme2}

Commençons par une proposition.

\begin{proposition2}
Soit $A$ un anneau muni d'une action d'un groupe fini $G$
d'ordre inversible sur $A$ et soit $B=\Fix_G(A)$ le sous-anneau
des invariants.
\begin{enumerate}
\item L'anneau $B$ est
\begin{enumerate}
\item nœthérien si $A$ l'est ;
\item réduit si $A$ l'est ;
\item local d'idéal maximal $𝔪 ∩ B$ si $A$ est local d'idéal
maximal $𝔪$, de corps résiduel isomorphe au sous-corps $\Fix_G(A ∕ 𝔪)$ de $κ=A ∕ 𝔪$.
\end{enumerate}
\item Le morphisme $\Spec(A) → \Spec(B)=\Spec(A) ∕ G$ est
\begin{enumerate}
\item fini si $A$ est nœthérien ;
\item génériquement étale si $A$ est de plus réduit.
\end{enumerate}
\end{enumerate}
\end{proposition2}

\begin{démo}
(i) Notons $\Tr$ le morphisme $B$-linéaire $\tr:A→  B$, $x\mapsto \frac{1}{|G|}∑_{g∈G}
g(x)$, parfois appelé « opérateur de Reynolds ».
Pour tout idéal $I$ de $B$, on a $IA ∩ B = I$. En effet,
l'inclusion $I ⊆ IA ∩ B$ est triviale et l'inclusion opposée
résulte du fait que si $x ∈ IA ∩ B$,
sa « trace » $x=\Tr(x)$ appartient, par $I$-linéarité, à $IB=I$.
On en déduit immédiatement l'énoncé a). L'énoncé b) est trivial.
Si $A$ est local, on a $A-𝔪=A^×$. Il résulte d'une part que $G$
stabilise globalement $𝔪$ et d'autre part que $\Fix_G(A) - \Fix_G(𝔪) = \Fix_G(A)^×$.
Ainsi, $B$ est maximal d'idéal $𝔫=\Fix_G(𝔪)$. Enfin, Le morphisme
canonique $B ∕ 𝔫 → \Fix_G(κ)$ déduit de l'inclusion canonique
$B ∕ 𝔫 → κ$ est un isomorphisme. En effet, si $a ∈A$ est un
relèvement arbitraire de $λ ∈ \Fix_G(κ)$, l'élément
$b=\Tr(a)$ en est un relèvement $G$-équivariant. Ceci achève
la démonstration du c).
(ii.a). Nous allons montrer que le morphisme entier $\Spec(A) →
\Spec(B)$ est fini par réduction au cas bien connu où $A$ est un corps.

❧ Réduction au cas réduit. Soient $N$ le nilradical de $A$ et $M=N∩B$ celui de $B$.
Pour chaque entier $i∈𝐍$, le $A/N$-module $N^i/N^{i+1}$ est de type
fini, car $A$ est supposé nœthérien, et nul pour $i ≫  0$.
Le module $\mathrm{gr}_{N}(A)=⊕_{n≥0}N^i/N^{i+1}$
est donc de type fini sur $\mathrm{gr}_{N}^0(A)=A/N$.
Si ce dernier est de type fini sur $B/M=\mathrm{gr}_{M}^0(B)$,
il en est de même de $\mathrm{gr}_{N}(A)$
sur $\mathrm{gr}_{M}(B)$ et finalement (\ac{iii}{2}{9}{, cor. 1})
de $A$ sur $B$, par complétude de l'anneau nœthérien $B$
pour la topologie $M$-adique.


❧ Réduction au cas d'un produit de corps.
Supposons $A$ réduit et considérons l'ensemble fini $\{𝔭_i\}_{i∈I}$
des idéaux premiers minimaux de $A$.
Pour chaque $i$, $𝔮_i=𝔭_i ∩ B$ est un idéal premier \emph{minimal} de $B$.
Cela résulte du théorème de Cohen-Seidenberg
(\ac{v}{2}{1}{, th. 1 et cor. 2}) et de
la transitivité de l'action de $G$ sur les fibres
de $\Spec(A) → \Spec(B)$ (\emph{op. cit.}, n°2, th. 2).
Soit $\Frac{ A}$ (resp. $\Frac{ B}$) l'anneau total des fractions de $A$ (resp. $B$) ;
c'est un produit de corps dans lequel $A$ (resp. $B$) s'injecte,
isomorphe au semi-localisé de $A$ en les $\{𝔭_i\}_{i∈I}$ (resp.
$\{𝔮_i\}_{i∈I}$). Soit $S=A-⋃_i 𝔭_i$ ; on a donc $\Frac{ A}=S^{-1}A$.
D'après (\emph{op. cit.}, §1, n°1, prop. 23), on a
$\Fix_G(S^{-1}A)=\Fix_G(S)^{-1}B$, de sorte que
$\Fix_G(\Frac{ A})=\Frac{ B}$ et $A⊗_{B} \Frac{ B}≅\Frac{ A}$.
Supposons $\Frac{\,A}$ fini sur $\Frac{\,B}$, de sorte qu'il existe d'après
l'isomorphisme précédent un nombre fini $n$ d'éléments de $A$ qui engendrent $\Frac{\,A}$
sur $\Frac{\,B}$. Observons que l'opérateur $\tr:A→B$ définit,
par composition avec le produit, un accouplement $A⊗_{B} A→ B$ 
qui est parfait sur les anneaux de fractions :
on se ramène à montrer que si $e_i$ est un idempotent
correspondant au facteur $K_i=\Frac{\,A/𝔭_i}$ de $\Frac{\,A}$, l'élément
$\tr(e_i)$ est non nul ;  il est en effet égal à $\frac{|G_i|}{|G|}$,
où $G_i$ est le stabilisateur de $e_i$.
Les $n$ éléments ci-dessus définissent donc un \emph{plongement} $A^G$-linéaire
de $A$ dans $B^n$. On peut conclure par nœthérianité.

❧ Réduction au cas d'un corps.
Soit donc $A=∏_i K_i$ un produit fini de corps et posons $X=\Spec(A)=∐_i η_i$.
Si $X=X₁∐X₂$, où $X₁$ et $X₂$ sont $G$-stables, $X/G=(X₁/G)∐(X₂/G)$ de sorte
que l'on se ramène immédiatement au cas où $X/G$ est connexe, c'est-à-dire où l'action
de $G$ est \emph{transitive}. Pour tout $i$, notons $G_i$ le groupe de
décomposition correspondant. D'après le cas classique (cas d'un corps),
$η_i → η_i/G_i$ est fini étale. Il en résulte que le morphisme
$X→ ∐ η_i / G_i$ est fini. Enfin, puisque pour tout $i$,
$η_i/G_i\iso X/G$ (\emph{loc. cit.}, §2, n°2, prop. 4), le résultat
(ii.a) en découle.

L'énoncé (ii.b) est désormais évident.
\end{démo}

\subsubsection{}Soient $A$ et $G$ comme dans l'énoncé
du théorème \ref{4-Cohen-Gabber équivariant}. Il résulte de la
proposition précédente l'on a l'égalité $\dim(B)=\dim(A)<+∞$, où
l'on note $B=\Fix_G(A)$. Nous noterons $d$ leur dimension commune.
Soit $B/I$ le quotient maximal $d$-équidimensionnel de $B$. D'après le théorème de
Cohen-Gabber \ref{4-Cohen-Gabber égale}, il existe un corps de représentants
$λ↪B/I$ et un système de paramètres $t₁,\dots,t_d$
de $B/I$ tel que $λ[[t₁,\dots,t_d]]→ B/I$ soit fini, génériquement étale.
On peut relever l'inclusion $λ↪B/I$ en une inclusion $λ↪B$ :
cela résulte par exemple, en caractéristique résiduelle positive
(seul cas non trivial), de la correspondance entre sous-corps de
représentants et relèvements d'une $p$-base donnée du corps résiduel.
Enfin, on peut relever le système de paramètres de $B/I$ en un système
de paramètres de $B$ : cela résulte, par dévissage, du lemme suivant.

\begin{lemme2}
Soient $A↠B$ une surjection d'anneaux locaux nœthériens et
$b∈\MM_B$ un élément \emph{sécant} pour $B$, c'est-à-dire
tel que $\dim(B/b)=\dim(B)-1$.
Il existe un relèvement de $b$ dans $A$ sécant pour $A$.
\end{lemme2}

Pour des généralités sur les suites sécantes,
voir par exemple \emph{op. cit.}, chap. VIII, §3. n°2.

\begin{démo}
On se ramène immédiatement au cas où $B=A/(f)$, $f∈A$.
Soit $a ∈ A$ un relèvement arbitraire de $b$ ;  par hypothèse, on a
$\dim(A/(f,a))=\dim(B)-1$.
Si $\dim(B)=\dim(A)-1$,
on a nécessairement $\dim(A/a)=\dim(A)-1$ car la dimension
chute d'au plus un par équation. Dans le cas contraire,
$f$ appartient à la réunion $⋃_{i=1}^{n}𝔭_i$, où les $𝔭_i$ sont les
idéaux premiers de $A$ de cohauteur
$\dim(A)$. Supposons que $f∈𝔭₁,\dots,𝔭_r$, et seulement
ces idéaux-ci. La conclusion ne peut être mise en défaut
que si $a+(f)⊆ ⋃_{i=1}^n 𝔭_i$, c'est-à-dire si tous les relèvements de $b$
sont non sécants. Pour chaque $i≤r$,
on a $a∉𝔭_i$ car $f$ appartenant à $𝔭_i$,
on aurait $\dim(A/(f,a))=\dim(A)$. Il en résulte notamment
que $r ≠ n$.
Il suffit donc de montrer que l'hypothèse
$a+(f) ⊆ ⋃_{i=r+1}^n 𝔭_i$ est absurde.
On aurait en effet $a+f^m=a+f ⋅ f^{m-1}∈  ⋃_{i=r+1}^n 𝔭_i$ pour tout $m$
et finalement $f^m(1-f^{m-m'})∈ 𝔭_i$ pour deux entiers $m>m'$ et un indice $r+1 ≤ i ≤ n$.
On en tire immédiatement $f∈𝔭_i$, ce qui est contraire à l'hypothèse.
\end{démo}

\subsubsection{}L'extension $κ/λ$ étant étale, car $λ=\Fix_G(κ)$,
le morphisme $κ→A/𝔪$ se relève \emph{uniquement} en un $λ$-homomorphisme $k→A$ ;
ce morphisme est $G$-équivariant. Le morphisme $A/B$ étant fini, génériquement étale,
ceci achève la démonstration du théorème \ref{4-Cohen-Gabber équivariant}.

\section{Autour du théorème de Epp}


\subsection{Énoncé (rappel)}

\subsubsection{}Si $X$ est un schéma \emph{réduit} n'ayant qu'un nombre fini
de composantes irréductibles, nous noterons $X_{\norm}$ son normalisé
(\ega{2}{6.3.6--8}).

\begin{theoreme2}[Helmut Epp, \cite{Eliminating@Epp}, théorème 1.9]\label{4-Epp}
Soit $T→S$ un morphisme local dominant de traits complets, de caractéristique
résiduelle $p>0$. Notons $κ_S$ et $κ_T$ leurs corps résiduels respectifs.
Supposons $κ_S$ \emph{parfait} et le sous-corps parfait maximal
de $κ_T$ \emph{algébrique} sur $κ_S$.
Il existe une extension finie de traits $S'→S$ telle que
le produit fibré réduit normalisé
$$
T':={(T\times_S S')_{\red}}_{\norm}
$$
ait une \emph{fibre spéciale réduite} au-dessus de $S'$.
\end{theoreme2}

\begin{remarque2}\label{4-remarque Epp}
En caractéristique mixte, le produit fibré $T\times_S S'$ 
est réduit. En effet, le morphisme $T'→S'$ (obtenu par
changement de base d'un plat) est plat, et $S'$ est
intègre si bien que l'anneau des fonctions de $T'$
s'injecte dans l'anneau des fonctions de sa fibre
générique. Il suffit donc de prouver que 
cette dernière est réduite. Or, en caractéristique nulle,
toute extension de corps est séparable. 
On vérifie
également sans difficulté que la conclusion du théorème est encore valable si 
l'on suppose seulement $S$ complet, mais pas
nécessairement $T$ (cf. \emph{loc. cit.}, §2).
\end{remarque2}

\subsection{Sorites}\label{4-stabilité E}

\subsubsection{}Nous dirons qu'une extension
de corps $K/k$ de caractéristique $p>0$ a la \emph{propriété de Epp} si tout
élément du sous-corps parfait maximal de $K$, $K^{p^∞}:=\cap_{i≥0} K^{p^{n}}$,
est algébrique séparable sur $k$. Pour $k$ parfait, c'est l'hypothèse
faite sur $κ_T/κ_S$ dans \ref{4-Epp}. Dans ce court
paragraphe, on rappelle quelques résultats élémentaires
de stabilité pour cette notion.
Commençons par un lemme.

\begin{lemme2}
Pour tout corps $K$ d'exposant caractéristique $p>1$, on a, dans une clôture séparable $K^{\sep}$ de $K$,
\[(K^{p^∞})^{\sep}=(K^{\sep})^{p^∞}.\]
\end{lemme2}

\begin{démo}
L'inclusion $(K^{p^{∞}})^{\sep}\subset (K^{\sep})^{p^∞}$ est
évidente : $K^{p^∞}$ est parfait donc toute extension algébrique,
en particulier sa clôture séparable
$(K^{p^{∞}})^{\sep}$, l'est également. Comme cette dernière est contenue
dans $K^{\sep}$, elle est également contenue dans son plus grand
sous-corps parfait $(K^{\sep})^{p^∞}$.


Réciproquement, considérons $x∈(K^{\sep})^{p^∞}$, et notons, pour chaque
entier $n≥0$, $x_n$ sa racine $p^n$-ième dans $K^{\sep}$ et $f_n$ son polynôme
minimal (\emph{unitaire}). Compte tenu d'une part de l'expression de $f_n$ en fonction
des polynômes symétriques en les conjugués galoisiens de $x_n$ et d'autre part
de l'injectivité et de l'additivité de l'élévation à la puissance $p^n$-ième,
on a l'égalité $f₀=f_n^{(p^n)}$, où $f_n^{(p^n)}$ est le polynôme
obtenu à partir de $f_n$ en élevant les coefficients à la puissance $p^n$-ième.
Il en résulte que les coefficients du polynôme minimal $f₀$ de $x$
appartiennent à $K^{p^∞}$.
\end{démo}

\begin{proposition2}[Cf. \cite{Eliminating@Epp}, §0.4]\label{4-prop Epp}
 Soit $k$ un corps d'exposant caractéristique $p$.
\begin{enumerate}
\item Soient $L/K$ et $K/k$ ayant la propriété
de Epp. Alors, $L/k$ a la propriété de Epp.
\item Toute extension finie de $k$ a la propriété de Epp.
\item Si $p>1$, pour tout entier naturel $d$, l'extension $\big(\Frac{\,k[[x₁,\dots,x_d]]}\big)/k$
a la propriété de Epp.
\item Si $p>1$, pour toute inclusion $k⊂A$,
où $A$ est un anneau local complet noethérien intègre,
induisant un isomorphisme sur les corps résiduels,
l'extension $(\Frac{\,A})/k$ a la propriété de Epp.
\end{enumerate}
\end{proposition2}

\begin{démo} Supposons immédiatement $p>1$ sans quoi
(i) et (ii) sont triviaux.\\
(i) Par hypothèse on a dans une clôture séparable de $L$ l'inclusion $L^{p^∞}\subset K^{\sep}$.
Comme le corps $L^{p^∞}$ est \emph{parfait}, on en déduit
que $L^{p^∞}\subset (K^{\sep})^{p^∞}=(K^{p^∞})^{\sep}\subset
k^{\sep}$, où l'égalité résulte du lemme précédent.

(ii) Toute extension étale a tautologiquement
la propriété de Epp. D'après (i), il reste à considérer
le cas d'une extension radicielle $K/k$.
Si elle est de hauteur $≤r$, on a $K^{p^r}\subset k$ et
en particulier $K^{p^∞}\subset k\subset k^{\sep}$.

(iii) Soit $A=k[[x₁,\dots,x_d]]$ et $K$ son corps
des fractions. Montrons que $K^{p^∞}=k^{p^∞}$.
Comme $K$ est contenu dans $k((x₁,\dots,x_{d-1}))((x_d))$,
on se ramène par récurrence au cas où $d=1$.
Tout élément non nul de $k((t))^{p^∞}$ a une valuation infiniment
$p$-divisible donc nulle, de sorte que
$k((t))^{p^∞}-\{0\}$ est contenu dans
$k[[t]]^×$ et finalement dans $k^{p^∞}$ par un calcul immédiat.
(iv) Cela résulte des observations précédentes et du théorème
de structure de·Cohen.
\end{démo}

\section{Le théorème de Cohen-Gabber en caractéristique mixte}

\subsection{Anneaux de Cohen et lissité formelle (rappels)}

\subsubsection{}Pour la commodité du lecteur, nous énonçons quelques résultats,
principalement dus à Cohen. Pour les démonstrations, nous renvoyons à \accourt{ix}{2} et \egazero{iv}{§19}.

\begin{définition2}[\egazero{iv}{19.3.1}]\label{4-19.3.1}
Soit $A$ un anneau topologique. Une $A$-algèbre topologique $B$
est dite \emph{formellement lisse} si pour toute $A$-algèbre topologique
\emph{discrète} $C$, et tout idéal \emph{nilpotent} $I$ de $C$,
tout $A$-morphisme continu $u:B→C/I$ se factorise en
$B\dessusdessous{v}{→}C\dessusdessous{φ}{↠}C/I$,
où $v$ est un $A$-morphisme continu et $φ$ l'homomorphisme canonique.
\end{définition2}

On dit aussi que $A→B$ est un \emph{morphisme formellement lisse}.
La proposition suivante énonce une propriété de relèvement
un peu plus générale que celle de la définition.

\begin{proposition2}[\emph{loc. cit.}, 19.6.1]\label{4-19.3.10}
Soient $A$ un anneau topologique, et $B$ un $A$-algèbre formellement lisse.
Soient $C$ une $A$-algèbre topologique, $I$ un idéal de $C$ vérifiant
les conditions suivantes :
\begin{enumerate}
\item $C$ est métrisable et complet ; 
\item $I$ est fermé et la suite $(I^n)_{n∈𝐍}$ tend vers zéro.
\end{enumerate}
Alors, tout $A$-morphisme continu $u:B→C/I$ se factorise en
$B\dessusdessous{v}{→}C↠C/I$,
où $v$ est un $A$-morphisme continu.
\end{proposition2}


Les deux théorèmes suivants donnent deux critères
importants de lissité formelle.

\begin{théorème2}[\emph{loc. cit.}, 19.6.1]
Une extension de corps munis de la topologie discrète
est formellement lisse si et seulement si l'extension
est \emph{séparable}.
\end{théorème2}

\begin{théorème2}[\emph{loc. cit.}, 19.7.1]\label{4-19.7.1}
Soient $A,B$ deux anneaux locaux nœthériens, $\got{m},\got{n}$
leurs idéaux maximaux respectifs et $k=A/\got{m}$ le corps
résiduel de $A$. Munissons $A$ et $B$ respectivement
des topologies $\got{m}$-adique et $\got{n}$-adique.
Soit $φ:A→B$ un morphisme \emph{local}, et posons $B₀=B⊗_A k$.
Les propriétés suivantes sont équivalentes :
\begin{enumerate}
\item $B$ est une $A$-algèbre formellement lisse ; 
\item $B$ est un $A$-module plat, et $B₀$ munie
de la topologie quotient est une $k$-algèbre formellement lisse.
\end{enumerate}
\end{théorème2}

Le théorème suivant, joint au précédent, est à la base
de la démonstration de l'existence des anneaux de Cohen
définis ci-après.
\begin{quote}
\emph{Dans les énoncés qui suivent, les anneaux locaux sont munis de la topologie 
de l'idéal maximal.}
\end{quote}

\begin{théorème2}[\emph{loc. cit.}, 19.7.2]\label{4-19.7.2}
Soient $A$ un anneau local nœthérien, $I$ un idéal strict, $A₀=A/I$, $B₀$
un anneau local nœthérien complet, $A₀→B₀$ un morphisme local formellement
lisse. Il existe alors un anneau local nœthérien complet $B$, un morphisme local
$A→B$ faisant de $B$ un $A$-module \emph{plat}, et un $A₀$-isomorphisme
$u:B⊗_A A₀\iso B₀$.
\end{théorème2}

\begin{définition2}[\emph{loc. cit.}, 19.8.4 et 5]\label{4-19.8.4}
\label{4-définition anneau Cohen}
On appelle \emph{anneau de Cohen} un anneau qui est soit
un corps de caractéristique nulle, soit un anneau de valuation discrète
complet, de corps résiduel de caractéristique $p>0$ et d'idéal
maximal engendré par $p$.
\end{définition2}

\begin{théorème2}[Cohen, \emph{loc. cit.}, 19.8.6 et 21.5.3]\label{4-19.8.6}
\begin{enumerate}
\item \label{1} Soient $W$ un anneau de Cohen de corps résiduel $K$, $C$ un anneau local nœthérien
complet, et $I$ un idéal strict de $C$. Alors, tout morphisme local
$u:W→C/I$ se factorise en $W\dessusdessous{v}{→}C→C/I$, où $v$ est local.
De plus, la factorisation est \emph{unique} si et seulement si 
$Ω¹_K=0$.
\item \label{2} Soit $K$ un \emph{corps}. Il existe un anneau de Cohen $W$ de corps 
résiduel isomorphe à $K$. Si $W′$ est un second anneau de Cohen, de corps 
résiduel $K′$, tout isomorphisme $u:K\iso K′$ provient par passage au
quotient d'un isomorphisme $v:W\iso W′$.
\end{enumerate}
\end{théorème2}

\subsubsection{}Rappelons que l'hypothèse $Ω¹_K=0$ est équivalente au fait
que $K$ est \emph{parfait} s'il est de caractéristique \mbox{$>0$} ou bien
est une extension algébrique de $𝐐$ s'il est de caractéristique nulle.

Signalons que si $K$ est \emph{parfait} de caractéristique $p>0$, le morphisme
$v$ de (\ref{2}) est \emph{unique}. Dans ce cas, $W$ est d'ailleurs 
isomorphe à l'anneau des vecteurs de Witt sur $K$.

\subsection{Le théorème de Cohen-Gabber en caractéristique mixte}

\subsubsection{}Soit $A$ un anneau local nœthérien complet de caractéristique résiduelle $p>0$.
Le schéma $X=\Spec(A)$ est de manière unique un $\Spec(𝐙_p)$-schéma. Notons
$X_p$ le sous-schéma fermé de $X$, fibre au-dessus du point fermé de $\Spec(𝐙_p)$.
Nous dirons qu'un ouvert $U⊂X$ est \emph{$p$-dense} si $U∩X_p$ est 
dense dans $X_p$.

\begin{théorème2}\label{4-Cohen-Gabber-Epp}
Soit $X=\Spec(A)$ un schéma local nœthérien complet normal de corps résiduel
$k$, de dimension $d≥2$ et de point générique de caractéristique nulle.
Il existe un morphisme fini surjectif $X'→X$, où $X'$
est normal intègre de corps résiduel $k'$, et
un morphisme fini surjectif $X'→\Spec(V[[t₁,\dots,t_{d-1}]])$,
où $V$ est un anneau de valuation discrète de corps résiduel $k'$,
\emph{étale} au-dessus d'un ouvert $p$-dense du but.
\end{théorème2}

La suite de ce paragraphe est consacrée à la démonstration
du théorème précédent.

\subsubsection{}Soit $X$ comme dans l'énoncé. Considérons 
le sous-corps parfait maximal $k₀=k^{p^∞}$ du corps résiduel $k$ de
$A=Γ(X,𝒪_X)$ et notons $W₀=W(k₀)$
l'anneau des vecteurs de Witt correspondant.
Il résulte du théorème de Cohen qu'il existe un unique morphisme
$X→S₀=\Spec(W₀)$ qui étende le morphisme 
$\Spec(k)→\Spec(k₀)$ entre les points fermés (\ref{4-19.8.6}, (i)).

Pour tout point maximal $\got{p}$ de la fibre spéciale $X_p$
de ce morphisme, l'anneau de valuation discrète $A_{\got{p}}$ a pour corps
résiduel $\Frac{\,A/\got{p}}$, où l'anneau $A/\got{p}$ est local
nœthérien complet intègre de corps résiduel $k$. D'après
\ref{4-prop Epp} (i) \& (iv), l'extension $\Frac(A/\got{p})/k₀$ a la propriété
de Epp. De tels idéaux $\got{p}$ étant
en nombre fini et la conclusion du théorème de Epp
(\ref{4-Epp}, \ref{4-remarque Epp}) étant
stable par changement de base fini car c'est un résultat
de lissité formelle, il existe donc un changement de base fini $S₀'=\Spec(W₀')→S₀$ 
tel que la fibre spéciale du produit fibré normalisé $X₀':=(X\times_{S₀}
S₀')_{\norm}=\Spec(A₀')$ soit réduite en ses points maximaux. (On utilise le fait que les points maximaux
de la fibre spéciale de $X₀'→ S₀'$ se trouvent au-dessus des points maximaux de la fibre
spéciale de $X→S₀$ ;  cf. p. ex. \egazero{iv}{16.1.6}.)

D'après le lemme suivant, la fibre spéciale du morphisme
$X₀'→ S₀'$ est alors réduite. 

\begin{lemme2}
Soit $X$ un schéma nœthérien normal. Tout diviseur de Cartier effectif 
génériquement réduit est réduit.
\end{lemme2}

\begin{démo}
On peut supposer $X$ affine égal à $\Spec(A)$ et le
diviseur de Cartier effectif défini par une
fonction $f∈A$. Soient $a∈A$ et $n≥1$ tels que 
$a^n∈fA$ ;  on souhaite montrer que $a∈fA$. 
L'anneau $A/f$ étant génériquement réduit,
l'élément $a/f$ de $\Frac{\,A}$ appartient à $A_𝔭$ 
pour tout idéal premier $𝔭$ contenant $f$ (un tel
idéal est de hauteur $1$ s'il est minimal
(\ac{viii}{3}{1}{\,prop. 1})). Il en est
évidemment de même pour $f∉𝔭$. L'anneau
$A$ étant normal, $⋂_𝔭 A_𝔭=A$ où $𝔭$ parcourt les idéaux de hauteur
$1$ (\emph{cf. loc. cit.} \textsc{vii}, §1, n°7, th. 4) 
de sorte que $a/f∈A$. CQFD.
\end{démo}

\subsubsection{}\label{4-fibration mixte} 
Notons $k₀'$ le corps résiduel de $W₀'$ 
, $\varpi'$ une uniformisante de $W₀'$, et considérons
une composante connexe $X'=\Spec(A')$ de $X₀'$ ; c'est un schéma fini
surjectif au-dessus de $X$. 
Soit $k'$ son corps résiduel.
L'inclusion $k₀'\inj k'$ déduite du morphisme $X'→S₀'$
est formellement lisse, car $k₀'$ est parfait, donc
se relève d'après \ref{4-19.7.1} et \ref{4-19.7.2} en
un morphisme \emph{formellement lisse} $W′₀→V$ où 
$V$ est un anneau local complet nœthérien. Cet anneau
est un anneau de valuation discrète.
L'anneau $A'/ϖ'$ étant \emph{réduit}, équidimensionnel de dimension $d-1$, de corps
résiduel $k'$, il existe d'après le théorème de
Cohen-Gabber (\ref{4-Cohen-Gabber égale}),
un relèvement $k₀'$-linéaire $k'↪ A'/ϖ'$ et 
des éléments $x_1,\dots,x_{d-1}$ dans l'idéal maximal de $A'/ϖ'$ 
tels que le morphisme induit
$k''[[t_1,\dots,t_{d-1}]]\ra A''/ϖ'$, envoyant
l'indéterminée $t_i$ sur $x_i$, soit
fini, \emph{génériquement étale} en haut et en bas.

Par lissité formelle de $W₀'→V$, le morphisme composé
$V→k'→A'/ϖ'$ se relève en un $W₀'$-morphisme $V→A'$.
En relevant les $x_i$ dans $A'$,
nous obtenons un morphisme $V[[t_1,\dots,t_{d-1}]]→A'$,
fini injectif (cf. p. ex. \egazero{iv}{19.8.8 (démonstration)}), 
\emph{étale} au-dessus du point générique de la fibre spéciale. CQFD.

\subsection{Le théorème de Cohen-Gabber premier à $ℓ$ en caractéristique mixte}

\begin{théorème2}\label{4-Cohen-Gabber-Epp équivariant}
Soit $X=\Spec(A)$ un schéma local nœthérien complet normal de
dimension $d≥2$, de corps résiduel $k$
de caractéristique $p>0$ et de point générique de caractéristique
nulle. Soit $ℓ$ un nombre premier différent de $p$.
Il existe alors :
\begin{enumerate}
\item un schéma local nœthérien intègre normal $Y$ muni d'une action
d'un $ℓ$-groupe fini $H$ et un morphisme fini surjectif
$H$-équivariant $Y → X$ tel que le quotient $Y ∕ H$ soit
de degré (générique) premier à $ℓ$ sur $X$ ;
\item un anneau de valuation discrète complet $V$ de même corps
résiduel $k ′$ que $Y$, de caractéristique mixte, muni d'une action de $H$
compatible avec son action sur $k ′$ ;
\item un morphisme local $Y→Y'=\Spec(V[[t_1,\dots,t_{d-1}]])$
qui soit fini, étale au-dessus d'un ouvert $p$-dense de $Y'$,
et $H$-équivariant avec action triviale de $H$ sur les $t_i$.
\end{enumerate}
\end{théorème2}

Ces morphismes sont représentés dans le diagramme ci-dessous,
où toutes les flèches sont des morphismes finis surjectifs.

$$
\xymatrix@C=3cm{
\Spec(V[[t_1,\dots,t_{d-1}]])=Y' & Y \ar[l]_(0.3){\text{$p$-génériquement étale}} \ar[d]
\ar[r] & Y/H \ar[dl]^{\textrm{ordre premier à } \ell} \\
& X & 
}
$$

\begin{remarque2}
Observons que les conditions (i) — (iii)
sur les morphismes $Y → X$ et $Y → Y ′$
n'entraînent pas que le schéma $Y/H$
soit étale au-dessus d'un ouvert $p$-dense de
$\Spec(\Fix_H(V)[[t_1,\dots,t_{d-1}]])=Y'/H$.
Voici un exemple, dû à Takeshi Saitô.
Soient $k$ un corps algébriquement clos de caractéristique $p>0$,
$W$ l'anneau des vecteurs de Witt sur $k$, $ℓ$ un nombre premier différent
de $p$, $A=W[[x, y]]/(x^ℓ y − p)$. Soient $W′=W[π]/(π^ℓ-p)$ et $B$ le
normalisé de $A⊗_W W′$, $W ′$-isomorphe à $W′[[x,z]]/(xz-π)$. Le groupe $H = μ_ℓ(k)$ agit sur $B$,
via son action sur $W′$ : $ζ\cdot x=x$ et $ζ\cdot z=ζz$.
Le morphisme $Y=\Spec(B)→X=\Spec(A)$ défini par $x↦x$ et $y↦z^ℓ$ satisfait les propriétés du théorème car
$Y/H⥲X$ et $Y→Y'=\Spec(W′[[x]])$, $x↦x$, est $p$-génériquement étale.
Cependant, $Y/H$ a une fibre spéciale isomorphe au schéma
non réduit $\Spec\big( k[[x,y]]/(x^ℓ\cdot y)\big)$.
Elle n'est donc pas étale au-dessus d'un
ouvert dense de la fibre spéciale de $Y ′$.
\end{remarque2}

La suite de ce paragraphe est consacrée à la démonstration du théorème
précédent. Notons que si la fibre spéciale $X_p$ de $X$ sur $\Spec(𝐙_p)$ est réduite,
ce théorème — comme le précédent — résulte simplement du
théorème \ref{4-Cohen-Gabber équivariant}, dans le cas particulier
où le groupe $G$ est trivial : on peut prendre $Y=X$ et $H$ trivial.

\subsubsection{}Considérons à nouveau le sous-corps parfait maximal $k₀$
du corps résiduel $k$ de $A$ et $W₀=W(k₀)↪A$ l'unique morphisme
relevant l'inclusion $k₀↪k$.
Soit $W₀^ν$ la clôture intégrale de $W₀$ dans $A$.

\begin{lemme2}
L'extension $W₀^ν/W₀$ est \emph{finie}, totalement ramifiée.
\end{lemme2}

\begin{démo}
Soit $W′/W₀$ une extension finie de traits, où $W′$ est
contenu dans $A$. Le corps résiduel de $W′$ est une
extension finie de $k₀$ ; c'est donc un corps parfait,
contenu dans $k$ et contenant $k₀$. Il est
donc égal à $k₀$ : l'extension est totalement ramifiée.
Le degré de l'extension $W′/W₀$ est par conséquent égal à son indice de
ramification, qui est \emph{majoré} par tout entier $N$ tel
que $p$ appartienne à $𝔪_A^N-𝔪_A^{N+1}$.
Si $W″$ est tel que le degré de l'extension $\Frac(W″) ∕ \Frac(W₀)$
soit maximal, on a nécessairement $W′ ⊆ W″$,
comme on le voit immédiatement en considérant
la sous-extension composée, dans $\Frac(A)$, des corps des fractions.
Ainsi, $W₀^ν=W″$ est fini sur $W₀$.
\end{démo}

\subsubsection{} D'après le théorème de Epp (\ref{4-Epp}),
il existe une extension finie d'anneaux de valuation
discrète $W₀^ν→W₀'$, que l'on peut supposer génériquement galoisienne
de groupe un groupe fini $G$, telle que la fibre spéciale sur $W₀'$ de la
normalisation $A'$ de $A⊗_{W₀^ν} W₀'$ soit \emph{réduite}. 
Observons que l'anneau $W₀^ν$ étant intégralement clos dans $A$,
l'anneau $A'$ est local.
Notons $k₀^ν$ (resp. $k₀'$) le corps résiduel de $W₀^ν$ 
(resp. $W₀'$) et $k'$ le corps résiduel de $A'$.
Choisissons des anneaux de Cohen $I(k')$ et $I({k'}^G)$ relatifs aux corps $k'$ et $k'^G$.
Il existe un morphisme $I({k'}^G)→I(k')$ relevant l'inclusion.
Ce morphisme étant fini étale entre anneaux locaux complets donc
henséliens, l'action du quotient $\Gal(k'/{k'}^G)$ de $G$ sur $k'$ 
se relève en une action $I({k'}^G)$-linéaire sur $I(k')$
(cf. p. ex. \cite{CL@Serre}, III, §5, th. 3).
Le corps $k₀'$ étant parfait, il existe d'après le théorème \ref{4-19.8.6}
un morphisme \emph{$G$-équivariant} $W(k₀')→I(k')$. 
Soient enfin $V=W₀'⊗_{W(k₀')} I(k')$, $ϖ$ une uniformisante de $W₀'$,
$\sur{A'}=A'/ϖA'$ et $H$ un $ℓ$-Sylow de $G$. D'après le théorème 
\ref{4-Cohen-Gabber équivariant}, il existe un morphisme \emph{fini, étale,
$H$-équivariant}, $φ:k'[[t₁,\dots,t_{d}]]→ \sur{A'}$, où les $t_i$
s'envoient dans $\sur{A'}^H$. Le morphisme ${A'}^H→ \sur{A'}^H$ étant
surjectif — comme cela se voit en utilisant la trace — on peut relever les images des $t_i$ 
en des $x'_i$ dans ${A'}^H$. De plus, par lissité formelle de $V/W₀'$, on peut relever
$k'→ \sur{A'}$ en un $W₀'$-morphisme $ψ:V→A'$ : cela résulte par exemple
de \egazero{iv}{19.3.10}. En procédant cran par cran,
et en considérant des isobarycentres dans les espaces affines
définis par le lemme bien connu suivant, on
constate qu'existe même un tel relèvement qui est \emph{$H$-invariant}.
\begin{lemme2}
Soit $(A,I)→ (B,J)$ un morphisme formellement lisse, où $A$ (resp. $B$) est
muni de la topologie définie par son idéal $I$ (resp. $J$). 
Soit $C↠C'$ une surjection de $A$-algèbres, de noyau $\mc{N}$ tel que
$\mc{N}²=I\mc{N}=0$. Alors, l'ensemble des $A$-linéaires relèvements d'un morphisme
$B→ C'$ à $C$ est soit vide soit un torseur sous $\mathrm{D\acute{e}r}_A (B,\mc{N})$.
\end{lemme2}

\subsubsection{}Le $X$-schéma $Y=\Spec(A')$ est bien fini $p$-génériquement étale 
sur $Y'=\Spec(V[[t₁,\dots,t_{d}]])$ si l'on envoie $V$ dans $A'$ par
$ψ$ comme ci-dessus et les variables $t_i$ sur les $x'_i$.
 Par construction $Y$ est, génériquement sur $X$, galoisien de groupe un sous-groupe de $G$ ;
son quotient $Y/H$ est donc génériquement d'ordre premier à $ℓ$.
Ceci achève la démonstration du théorème.

\ifx\danslelivre\undefined
\bibliography{../fichiers-configuration-latex/bibliographie}
\bibliographystyle{../fichiers-configuration-latex/nomcomplet-date}
\end{document}
\fi
